\section{Samples and observations}

A dataset is usually written as a finite list
\[
x_1,x_2,\ldots,x_n.
\]
The number \(n\) is called the sample size. Each \(x_i\) is one observed value.
For example, if five waiting times are recorded, then a dataset might be
\[
2.1,\quad 3.4,\quad 1.8,\quad 4.0,\quad 2.7.
\]
Here \(n=5\), and each number is one observation.

\begin{definitionbox}
A \textbf{sample} is a finite collection of observed values
\[
x_1,x_2,\ldots,x_n.
\]
The values \(x_i\) are called \textbf{observations}, and \(n\) is the
\textbf{sample size}.
\end{definitionbox}

This definition is deliberately simple. At the descriptive level, the data are
already observed. We can sort them, count them, average them, plot them, and
summarize them.

\subsection*{Random variables before observation, numbers after observation}

In probability modelling, before data are observed, we often write
\[
X_1,X_2,\ldots,X_n.
\]
These symbols represent random variables. After observation, their realized
values are written as
\[
x_1,x_2,\ldots,x_n.
\]
This distinction is small in notation but large in meaning.

\begin{center}
\begin{tabular}{lll}
\toprule
Notation & Meaning & Level \\
\midrule
\(X_i\) & random variable before observation & model \\
\(x_i\) & observed numerical value & data \\
\bottomrule
\end{tabular}
\end{center}

For instance, before measuring tomorrow's temperature at noon, the temperature
can be modelled as a random variable \(X\). After the measurement is made, the
recorded value, say \(18.6\), is a number. The random variable describes
uncertainty before observation; the observed value is what actually appeared.

\begin{remarkbox}
The same letter in uppercase and lowercase often marks this difference:
uppercase for the random variable, lowercase for the observed value.
\end{remarkbox}

\subsection*{Samples from a population}

In applications, data are often interpreted as observations from a larger
population. The population may be a real finite population, such as all students
in a university, or an idealized population represented by a probability model.

For example:
\begin{itemize}
    \item a sample of heights may be taken from a population of people;
    \item a sample of delivery times may be taken from a logistics process;
    \item a sample of simulated values may be generated from a normal
    distribution;
    \item a sample of measurement errors may be collected from repeated
    measurements of the same quantity.
\end{itemize}

In each case, the observed data are finite. The population or probability model
is usually larger than the sample and may not be fully known.

\subsection*{Independent and identically distributed observations}

A common modelling assumption is that observations arise from random variables
\[
X_1,X_2,
\ldots,X_n
\]
that are independent and have the same probability law. This is often abbreviated
as \textbf{i.i.d.}, meaning \textbf{independent and identically distributed}.

This assumption means two things:
\begin{itemize}
    \item the observations are produced by the same probabilistic mechanism;
    \item knowing one observation does not change the probability law of the
    others.
\end{itemize}

The i.i.d. assumption is powerful, but it is an assumption. It should not be
silently imposed on every dataset. Measurements over time may be dependent.
People in the same household may be similar. Repeated measurements may be
affected by the same instrument error. The model must be checked against the
context.

\subsection*{Real data and computer-generated data}

This chapter uses both real-data language and computer-generated samples. A
computer-generated sample is not real-world evidence by itself, but it is useful
for learning. It allows us to see how definitions behave when the underlying
model is known.

For example, a computer may generate values from a uniform distribution on
\([0,1]\). Then we know the theoretical CDF is
\[
F(x)=x,
\qquad 0\leq x\leq1.
\]
After generating a finite sample, we can compute the empirical CDF and compare it
with the theoretical one. This comparison helps us understand the relationship
between a model and a finite sample from that model.

\begin{summarybox}
At the model level, we work with random variables \(X_1,\ldots,X_n\). At the data
level, we work with observed values \(x_1,\ldots,x_n\). Descriptive statistics are
computed from the observed values, while probability laws describe the random
mechanism that may have produced them.
\end{summarybox}
